\documentclass[12pt]{article}
\usepackage{lmodern}
\usepackage{xcolor}
\usepackage{amsmath, amsthm, amssymb}
\usepackage[margin=1in]{geometry}
\input{_macros}

\newtheorem{theorem}{Theorem}
\newtheorem{lemma}{Lemma}
\newtheorem{assumption}{Assumption}
\newtheorem{definition}{Definition}

\begin{document}

\section*{Section 5: Synthesis --- Proof of the Main Theorem}

This section assembles the four independently proved conclusions
C1--C4 into a proof of the Main Theorem (Theorem~1 of the problem statement).

\subsection*{Recalled Conclusions}

\begin{itemize}
  \item \textbf{C1 (Capability Inevitability, proved in Section 1):}
    Under A1, the scaling-law AI system exhibits exponentially decreasing
    loss, accelerating capability growth ($\kappa(t) = e^{\beta t}/A$
    is strictly convex), and economically dominates the pre-AI baseline
    in every period.

  \item \textbf{C2 (Economic Lock-In, proved in Section 2):}
    Under A2 and A4, for any organization at integration depth $d \geq 1$
    and any discount rate $\rho > 0$:
    \[
      V_{\mathrm{rev}}(d,\rho) - V_{\mathrm{cont}}(d,\rho)
      \;\leq\; -c\,d^2 - \frac{\gamma\Pi_0}{\rho} \;<\; 0.
    \]
    Reversion is never individually economically rational; the case
    against reversion strengthens monotonically with depth.

  \item \textbf{C3 (Scientific Force Multiplication, proved in Section 3):}
    Under A5, AI has made beyond-record contributions in at least two
    independent scientific fields, documented a $>2\times$ time speedup
    and a $>100\times$ discovery-scale ratio, spanning four independent
    scientific disciplines.

  \item \textbf{C4 (Structural Unavoidability, proved in Section 4):}
    Under A3, A4, and A4b, AI is embedded in four critical societal
    systems with switching costs $S(d_X) \geq B_X$ (financially
    prohibitive); continued use only increases these costs; no viable
    non-AI substitute exists; AI's role is structurally unavoidable.
\end{itemize}

\subsection*{Main Theorem}

\begin{theorem}[AI Is the Future --- Main Result]
  Under Assumptions A1, A2, A3, A4, A4b, and A5, Artificial Intelligence
  constitutes the defining technological paradigm of humanity's future
  development, as evidenced jointly by C1--C4:

  \begin{enumerate}
    \item \emph{Technological Inevitability (C1):} AI capability growth
      is self-reinforcing, accelerating, and makes the pre-AI baseline
      economically dominated in every period.

    \item \emph{Economic Lock-In (C2):} At any integration depth
      $d \geq 1$, AI adoption is individually economically irreversible:
      no positive discount rate makes reversion rational.

    \item \emph{Scientific Force Multiplication (C3):} AI is a
      cross-domain scientific accelerant with demonstrated beyond-record
      contributions, quantified speedups, and global disciplinary reach.

    \item \emph{Structural Unavoidability (C4):} AI is already embedded
      in critical societal infrastructure at financially prohibitive
      switching-cost depth, with no viable substitute, making its
      continued role structurally unavoidable.
  \end{enumerate}

  Together, C1--C4 establish a convergent, mutually reinforcing case:
  AI is not a transient technology but the defining paradigm of
  humanity's future.
\end{theorem}

\begin{proof}
  Each of C1, C2, C3, C4 is an independently proved conclusion:
  \begin{itemize}
    \item C1 is proved in Section~1 from A1 alone (Theorem~1 of Section~1).
    \item C2 is proved in Section~2 from A2 and A4 alone (Theorem~1 of Section~2).
    \item C3 is proved in Section~3 from A5 alone (Theorem~1 of Section~3).
    \item C4 is proved in Section~4 from A3, A4, and A4b (Theorem~1 of Section~4).
  \end{itemize}
  The dependency graph is acyclic: no conclusion is used to prove any
  other.  All assumptions A1, A2, A3, A4, A4b, A5 are empirically
  documented and declared in the problem statement.

  The four pillars are \emph{convergent}: C1 shows AI's capability will
  keep improving (no technical plateau), C2 shows existing adopters
  cannot economically exit (lock-in), C3 shows AI accelerates the
  frontier of human knowledge (scientific compounding), and C4 shows
  AI is already embedded at the infrastructure level with no exit path.
  No single counter-argument can undercut all four simultaneously.

  Therefore, under the stated assumptions, AI constitutes the defining
  technological paradigm of humanity's future development.  \qed
\end{proof}

\end{document}
